\section*{Hoofdstuk \ref{ch:g9:elements-universe-earth} --- Elementen in het heelal en op aarde}

\begin{solution}{exo:g9:elements-universe-earth:1}
Waterstof (ongeveer \qty{74}{\percent}) en helium (ongeveer
\qty{25}{\percent}).
\end{solution}

\begin{solution}{exo:g9:elements-universe-earth:2}
In beide zuurstof.
\end{solution}

\begin{solution}{exo:g9:elements-universe-earth:3}
Binnen in sterren, en bij hun ontploffingen.
\end{solution}

\begin{solution}{exo:g9:elements-universe-earth:4}
Ongeveer \qty{3.5}{\percent}.
\end{solution}

\begin{solution}{exo:g9:elements-universe-earth:5}
$46.6 + 27.7 = \qty{74.3}{\percent}$: ongeveer drie kwart.
\end{solution}

\begin{solution}{exo:g9:elements-universe-earth:6}
Het zijn heel lichte gassen: ze zijn uit de jonge aarde de ruimte in
ontsnapt (en helium vormt geen verbindingen die het in de gesteenten
hadden kunnen vasthouden).
\end{solution}

\begin{solution}{exo:g9:elements-universe-earth:7}
$0.23 \times 70 = \qty{16.1}{kg}$.
\end{solution}

\begin{solution}{exo:g9:elements-universe-earth:8}
Zuurstofatomen zijn 16 keer zwaarder dan waterstofatomen: minder
zuurstofatomen wegen toch meer.
\end{solution}

\begin{solution}{exo:g9:elements-universe-earth:9}
Ongeveer $0.081 \times 1000 = \qty{81}{kg}$ aluminium.
\end{solution}

\begin{solution}{exo:g9:elements-universe-earth:10}
Zuurstof en calcium (en, buiten de acht belangrijkste elementen van de
aardkorst, komen waterstof, koolstof en fosfor in kleine hoeveelheden in
de aardkorst voor).
\end{solution}

\begin{solution}{exo:g9:elements-universe-earth:11}
$40 \times 1.67 \times 10^{-27} = \qty{6.68e-26}{kg}$;
$1.06 / 6.68 \times 10^{-26} \approx 1.6 \times 10^{25}$ calciumatomen.
\end{solution}

\begin{solution}{exo:g9:elements-universe-earth:12}
\omterm{def:g7:atoms-and-molecules:atom}{Atomen} worden door \omterm{def:g5:irreversible-changes:chemical-change}{chemische veranderingen} niet gemaakt en niet
vernietigd: ze gaan alleen over van de ene verbinding naar de andere
(gesteente, gas, suiker, levend lichaam). De koolstofatomen van de aarde
worden steeds opnieuw gebruikt.
\end{solution}

\begin{solution}{pb:g9:elements-universe-earth:1}
\textbf{1.} $0.61 \times 70 = \qty{42.7}{kg}$.

\textbf{2.} Koolstof $0.23 \times 70 = \qty{16.1}{kg}$; waterstof
$0.10 \times 70 = \qty{7.0}{kg}$.

\textbf{3.} $61 + 23 + 10 = \qty{94}{\percent}$.

\textbf{4.} Water, \ce{H2O}.

\textbf{5.} $16 \times 1.67 \times 10^{-27} = \qty{2.67e-26}{kg}$.

\textbf{6.} Koolstof: $12 \times 1.67 \times 10^{-27} = \qty{2.00e-26}{kg}$;
waterstof: \qty{1.67e-27}{kg}.

\textbf{7.} 16.

\textbf{8.} Een \omterm{def:g9:inside-the-atom:nucleus}{elektron} is ongeveer 1836 keer lichter dan een \omterm{def:g9:inside-the-atom:proton}{nucleon}:
de \omterm{def:g9:inside-the-atom:nucleus}{elektronen} zijn minder dan een duizendste van de massa.

\textbf{9.} $7.0 / 1.67 \times 10^{-27} \approx 4.2 \times 10^{27}$
waterstofatomen.

\textbf{10.} $16.1 / 2.00 \times 10^{-26} \approx 8.0 \times 10^{26}$
koolstofatomen.

\textbf{11.} $42.7 / 2.67 \times 10^{-26} \approx 1.60 \times 10^{27}$
zuurstofatomen.

\textbf{12.} Onze telling, $1.60 \times 10^{27}$, klopt met de schatting
$1.61 \times 10^{27}$: een lichaam van \qty{70}{kg} bevat ongeveer $1.6 \times
10^{27}$ zuurstofatomen.
\end{solution}
